\documentclass{article}
\usepackage[preprint]{hackathon_style}

\usepackage[utf8]{inputenc}
\usepackage[T1]{fontenc}
\usepackage{hyperref}
\usepackage{url}
\usepackage{booktabs}
\usepackage{amsmath}
\usepackage{amsfonts}
\usepackage{microtype}
\usepackage{listings}

\hypersetup{colorlinks=true, linkcolor=black, citecolor=black, urlcolor=blue}
\lstset{basicstyle=\ttfamily\scriptsize, breaklines=true, frame=single,
        framesep=4pt, rulecolor=\color{gray!50}, columns=fullflexible,
        keepspaces=true, showstringspaces=false}

\title{\textsc{Endogrid}: Measuring the Cost of Reasoning\\on a Long-Horizon
       Economic Planning Benchmark}

\author{%
  \normalsize
  Hackathon submission\\
  \texttt{worldforge\_bench}
}

\begin{document}
\maketitle

\begin{abstract}
Architectures such as the Dragon Hatchling \citep{kosowski2025bdh} and BDH-CQ
\citep{engdahl2026bdhcq} are motivated by a specific bet: that adaptation
through \emph{recurrent state} can match explicit reasoning at a fraction of the
inference cost. Testing that bet needs environments where both routes can be
run under one budget and one score. We present \textsc{Endogrid}, a
deterministic ten-simulated-year infrastructure-development benchmark whose
objective is the equity value of a business rather than any quantity the agent
can monotonically accumulate, and we use it to measure the cost of reasoning
directly. Under a matched twenty-decision budget on the same map, a
state-based policy that performs no search, no backtracking and no weight
updates scores $40.8$; the best language-model agent scores $43.4$. The policy
reaches $94\%$ of that score at $1/670$ of the cost --- $32{,}870$ points per
dollar against $53.9$. Meanwhile an agent restricted to a fixed
terrain$\rightarrow$machine lookup table scores $6.8$, \emph{below} the $16.3$
of doing nothing at all, confirming the environment penalises memorised
surface correlations rather than rewarding them. We are explicit about what
this is and is not: we did not run a BDH model. What we contribute is the
harness, the protocol and the measured frontier on which that comparison can be
made, together with the finding that on this task the expensive route buys
little.
\end{abstract}

\section{The question}

The interesting claim shared by the Dragon Hatchling \citep{kosowski2025bdh} and
BDH-CQ \citep{engdahl2026bdhcq} is not that recurrent-state models are more
accurate. It is that a large part of what explicit reasoning buys can be got
more cheaply. BDH is evaluated on Sudoku-Extreme at $97.4\%$ top-1 \emph{with no
chain-of-thought and no backtracking}; BDH-CQ adapts to an unseen ARC task
through recurrent state alone, with no parameter updates at inference. Both are
statements about the price of correctness, not only its ceiling.

That claim is hard to test on puzzle benchmarks, because a puzzle is scored
right or wrong and the cost axis is left implicit. What is missing is an
environment where (i) an agent must plan over a long horizon in an unfamiliar
world, (ii) the objective cannot be climbed by doing more of the rewarded thing,
and (iii) both a cheap state-based policy and an expensive reasoning system can
be run under one budget and scored on one number.

We built that environment and measured the frontier.

\paragraph{What we claim, and what we do not.} We did not run a BDH model.
Every result here is a scripted policy or a language-model agent. The
contribution is the harness and the measurement: a benchmark where the
cost of reasoning is visible, a protocol borrowed from the BDH and ARC-AGI-2
evaluations, and the empirical observation that on this task a policy with no
search loop is within $6\%$ of the best language model at three orders of
magnitude less cost. That is a motivation for architectures like BDH, not
evidence for them.

\section{The environment}

An agent is given a procedurally generated terrain, a starting balance sheet of
\$25M, and ten simulated years. It sites solar, wind and hydro generators,
orients them, wires them into demand zones, finances them with cash and debt,
and may hedge output with long-term contracts. Each simulated hour the physical
layer resolves --- sun position, wind field, water cycle, river flow, machine
output, transmission --- and then a market clears.

\subsection{Why the objective is not monotone}

Most agent benchmarks score an accumulable quantity, so an agent climbs by doing
more of the rewarded thing. Here the price of electricity is endogenous: it
clears every tick against a residual demand curve served by a background thermal
fleet in merit order. Renewables bid at zero marginal cost, so the agent's own
output displaces the marginal unit and lowers the price it earns:

\begin{center}
\small
\begin{tabular}{rrl}
\toprule
Agent supply (kW) & Clearing price (\$/MWh) & Marginal unit \\
\midrule
0        & 79.10    & CCGT \\
10{,}000 & 75.25    & lignite \\
20{,}000 & 9.00     & nuclear \\
30{,}000 & $-3.12$  & oversupply \\
60{,}000 & $-25.00$ & price floor \\
\bottomrule
\end{tabular}
\end{center}

Each asset is discounted at its \emph{own} realised capture price, so an agent
that floods the market marks down the assets it already owns. Building more can
be worth less, and the naive maximising policy defeats itself. Wealth can fall:
capital is scarce, leverage is priced, covenants bind, and insolvency ends the
episode.

\subsection{Rules that must be inferred, not memorised}

No terrain carries an energy multiplier. ``Sand means solar'' and ``mountain
means wind'' are both false by construction --- a leeward mountain cell has
\emph{worse} wind than flat ground, because of an orographic wind-shadow term.
Output traces through chained physical variables: irradiance through a
terrain-obstruction ray-cast and a real incidence cosine against panel tilt and
azimuth; turbine power through $\tfrac{1}{2}\rho A v^3$ with a
$\cos(\text{yaw})$ misalignment term that is a hard zero at $90^\circ$; hydro
through $\rho g Q H$ with floors on discharge and head.

This is a deliberate defence against causal confusion in the sense of
\citet{tien2023causal}: an agent that latches onto a correlated surface feature
instead of the causal chain should be punished. Section~\ref{sec:results} shows
it is.

\section{Protocol}
\label{sec:protocol}

We borrow the evaluation logic of the BDH line of work, and say plainly that it
is the logic we are borrowing, not the model.

\paragraph{No weight updates during the run.} An agent plays ten simulated years
with no retraining, no fine-tuning and no gradient step. Only in-episode state
changes: cash, built assets, and whatever the agent has learned about
\emph{this} map's wind, water and terrain. This mirrors BDH-CQ's adaptation
through recurrent state rather than parameters \citep{engdahl2026bdhcq}. It is a
meaningful constraint rather than an arbitrary one precisely because the
alternative exists: the Hierarchical Reasoning Model \citep{wang2025hrm} takes a
backward pass per unseen task, and our agents are given no such route. Fixing
the constraint makes the two families comparable.

\paragraph{Memory is environmental, not a cache.} The recurrent state an agent
must track is physical: soil moisture that carries rainfall forward, river
discharge that lags its catchment, accumulated evidence about a prevailing wind
bearing that mean-reverts seasonally, machine health that degrades. These are
closer to the fast-weight, Hebbian synaptic memory of \citet{kosowski2025bdh}
than to a key--value cache of past tokens: the state \emph{is} the world's, it
persists whether or not the agent attends to it, and it is what makes a
consequence delayed.

\paragraph{Correctness without an explicit search loop.} BDH is tested on
whether it reaches correct Sudoku solutions with no chain-of-thought and no
backtracking \citep{kosowski2025bdh}. We apply the same behavioural signature.
Our \texttt{heuristic} agent performs no search: it reads the fields at each
candidate cell, ranks by expected return on capital, commits, and never revises
by exploring alternatives. It cannot backtrack --- a built machine is a spent
\$1.45M. Whether such a policy reaches good placements is the same question BDH
is asked.

\paragraph{Held-out seeds.} Following the private-set methodology of
ARC-AGI-2 \citep{chollet2025arcagi2}, physical constants are fixed while map
layout, river paths, ridge geometry and demand placement vary with the seed.
Seeds 1--8 are development; \textbf{seeds 9 and 10 were not run until final
scoring}. A generalisation gap is reported rather than assumed.

\begin{table}[t]
\centering\small
\setlength{\tabcolsep}{4pt}
\caption{Representative evaluation efforts, and where \textsc{Endogrid} sits.
\emph{Adaptation} is the route an agent has to a task it has not seen.
\emph{Cost axis} asks whether the benchmark scores what correctness cost, not
only whether it was reached.}
\label{tab:related}
\begin{tabular}{llllll}
\toprule
& World & Objective & Horizon & Adaptation & Cost axis \\
\midrule
Sudoku-Extreme \citep{kosowski2025bdh} & fixed puzzle & correct grid & one shot & recurrent state & no \\
ARC-AGI-2 \citep{chollet2025arcagi2} & held-out tasks & correct grid & one shot & in-context & no \\
HRM \citep{wang2025hrm} & puzzle suite & correct grid & one shot & backward pass & params only \\
TRM \citep{jolicoeurmartineau2025trm} & puzzle suite & correct grid & one shot & recursive & params only \\
Crafter \citep{hafner2022crafter} & procedural grid & achievements & episode & policy weights & no \\
MineDojo \citep{fan2022minedojo} & open world & task suite & episode & policy weights & no \\
\midrule
\textbf{\textsc{Endogrid}} & procedural grid & \textbf{equity value} & \textbf{10 sim-years} & in-episode state & \textbf{dollars} \\
\bottomrule
\end{tabular}
\end{table}

\paragraph{One world, many tiers.} That a single procedurally generated grid
world with graded difficulty is a legitimate evaluation design rather than a toy
was established by Crafter \citep{hafner2022crafter}, and extended to
open-ended task suites by MineDojo \citep{fan2022minedojo}. Our capability
targets --- spatial reasoning over terrain and wind geometry, construction
planning, memory of prior environmental outcomes, causal rule learning, and
transfer to unfamiliar layouts --- follow that lineage.

\section{Agents}

\begin{center}
\small
\begin{tabular}{lp{9.4cm}}
\toprule
Agent & What it does \\
\midrule
\texttt{donothing} & Holds its starting capital. The floor any agent must clear
to have created value. \\
\texttt{lookup} & Consults a fixed terrain$\rightarrow$machine table at a fixed
orientation. Never observes wind, sun, price or its own books. The
memorisation control. \\
\texttt{random} & Builds legal machines at random cells and orientations. The
noise floor. \\
\texttt{heuristic} & Ranks sites by the fields actually present at each cell,
orients to the resource, wires to the nearest reachable zone, and stops building
once its own supply would depress the price. No search, no backtracking. \\
\midrule
LLM agents & Opus 5, Sonnet 5, Haiku 4.5, playing blind through a command
interface. Given the action schema and the objective only --- no description of
the price mechanism, of what drives output, or of how assets are valued --- and
instructed not to read the implementation. \\
\bottomrule
\end{tabular}
\end{center}

% Generated from runs/ by make_results.py, so the paper cannot drift from
% what the code produced. Absent until the benchmark has been run.
\IfFileExists{results.tex}{\section{Results}
\label{sec:results}

\begin{table}[t]
\centering\small
\caption{Ten seeds, ten simulated years each, matched twenty-decision budget.
Equity is terminal equity value in millions against a \$25M start; GWh is energy
delivered; LCOE is levelised cost in \$/MWh; \emph{Ins.} counts insolvencies.}
\label{tab:multi}
\begin{tabular}{lrrrrrrrrr}
\toprule
Agent & Score & sd & min & max & Equity & Return & GWh & LCOE & Ins. \\
\midrule
\texttt{heuristic} & 43.8 & 11.5 & 14.6 & 55.1 & 46.5 & +86\% & 381 & 44 & 0 \\
\texttt{random} & 22.2 & 11.4 & 10.7 & 43.8 & 28.3 & +13\% & 52 & 157 & 0 \\
\texttt{hebbian-mwh} (energy target) & 20.9 & 6.9 & 14.1 & 34.9 & 22.6 & $-$10\% & 436 & 60 & 2 \\
\texttt{hebbian} (BDH-inspired) & 18.2 & 3.3 & 14.7 & 25.5 & 22.9 & $-$9\% & 292 & 71 & 2 \\
\texttt{donothing} & 16.3 & 0.0 & 16.3 & 16.3 & 25.0 & +0\% & 0 & 0 & 0 \\
\texttt{lookup} & 9.1 & 2.5 & 4.4 & 12.2 & 17.0 & $-$32\% & 86 & 290 & 0 \\
\bottomrule
\end{tabular}
\end{table}

\paragraph{Memorisation loses money.} \texttt{lookup} scores 9.1, below the 16.3 of holding capital and building nothing, with a mean return of -32\%. An agent applying a fixed terrain$\rightarrow$machine table does not merely fail to excel; it destroys capital, because it sites without regard to the resource that actually reaches a cell and orients every machine identically regardless of the wind. This is the falsifiable form of the anti-memorisation claim: if a lookup table ever became competitive, the environment would be rewarding the surface correlations it was built to punish \citep{tien2023causal}.

\paragraph{The BDH-inspired agent fails, and how it fails is the point.} \texttt{hebbian} carries a fast-weight associative memory over binned site observations, zero-initialised every episode, updated by a local gradient-free Hebbian rule from realised outcomes, with no search and no backtracking. It scores 18.2 against \texttt{heuristic}'s 43.8, returns -9\\%, and goes insolvent on 2 of ten seeds. Redirecting its reward from delivered energy to realised value barely moved it (20.9 to 18.2) and did not curb the overbuilding: 52 machines on average against \texttt{heuristic}'s 16.

The diagnosis is structural rather than a tuning failure. A per-site associative memory maps \emph{features of a cell} to outcomes, but cannibalisation is not a property of any cell --- ``the next machine earns less because I already built ten'' is a property of total portfolio supply. No weight over wind, irradiance, head or stability can express it, so the agent keeps finding sites its memory rates highly and keeps depressing the price it earns on all of them. This is a concrete statement of what state a recurrent architecture would need to carry on this benchmark: not richer features per site, but a representation of its own aggregate position. We report it as a negative result. It is the most useful thing the benchmark told us.

\paragraph{Held-out seeds.} \texttt{heuristic} averages 46.0 on the eight development seeds and 34.9 on seeds 9 and 10, which were not run until final scoring --- a relative gap of +24.2\%, within noise, so we see no evidence of layout fitting. The policy has no parameters to fit, so this is a check that the \emph{environment} does not vary wildly by layout, as much as a check on the agent.

\begin{table}[t]
\centering\small
\caption{Cost of reasoning. All agents on seed 1 under the same twenty-decision
budget. Language-model cost is the whole-episode token total at published
first-party rates; scripted cost prices CPU at \$0.10/hour, a deliberately
generous figure. \emph{Score/\$} is the resulting accuracy-per-dollar. Scripted rows come from
the same engine version as Table~\ref{tab:multi}. Haiku~4.5's row predates a
correction to hydro siting and is contaminated by it; the other two
language-model rows built no affected assets and are unaffected.}
\label{tab:cost}
\begin{tabular}{lrrrrr}
\toprule
Agent & Score & Equity & Cost & Minutes & Score/\$ \\
\midrule
Sonnet 5 & 43.4 & 46.8 & \$0.80 & 3.8 & 54 \\
Opus 5 & 33.9 & 37.9 & \$2.26 & 6.3 & 15 \\
Haiku 4.5 & 8.2 & 18.7 & \$0.35 & 2.1 & 24 \\
\midrule
\texttt{heuristic} & 40.8 & 44.9 & \$0.0011 & 0.7 & 36,960 \\
\texttt{random} & 43.8 & 44.9 & \$0.0010 & 0.6 & 43,050 \\
\texttt{hebbian-mwh} (energy target) & 34.9 & 44.7 & \$0.0009 & 0.5 & 39,665 \\
\texttt{hebbian} (BDH-inspired) & 15.9 & 25.7 & \$0.0010 & 0.6 & 16,264 \\
\texttt{donothing} & 16.3 & 25.0 & \$0.0007 & 0.4 & 22,149 \\
\texttt{lookup} & 11.7 & 22.8 & \$0.0011 & 0.7 & 10,572 \\
\bottomrule
\end{tabular}
\end{table}

Table~\ref{tab:cost} is the result this paper exists for. On the same map, under the same budget, Sonnet 5 scores 43.4 and \texttt{heuristic} scores 40.8. The language model is ahead --- we are not claiming otherwise. But it costs \$0.80 and 3.8 minutes to get there, against a tenth of a cent and 0.7 minutes. The policy with no search loop reaches 94\% of the score at roughly $1/686$ of the cost, or 686$\times$ the accuracy per dollar.

\paragraph{Why this is the interesting axis.} The same shape of argument runs through the recursive-reasoning literature: TRM reaches competitive accuracy with 7M parameters where HRM uses 27M \citep{jolicoeurmartineau2025trm,wang2025hrm}, and BDH-CQ adapts through recurrent state rather than a backward pass per task \citep{engdahl2026bdhcq}. What those results assert about parameter counts, Table~\ref{tab:cost} measures in dollars on a task with a long horizon and an economic objective. A model that matched the language-model score at the state-based policy's cost would dominate this frontier outright. Nothing we ran occupies that corner --- which is precisely the corner BDH-style architectures claim.

}{%
  \section{Results}
  \label{sec:results}
  \emph{Results not yet generated. Run the benchmark and the generator:}
  \begin{lstlisting}[language=bash]
wfbench bench --agents donothing,lookup,random,heuristic \
  --seeds 1-10 --years 10 --interval 4320 --out runs/hackathon_10seeds.json
python3 paper-hackathon/make_results.py
  \end{lstlisting}
}

\section{Limitations}

\textbf{No BDH model was run.} This is a harness and a measurement, not an
architecture comparison. The frontier we report is where a BDH-style model would
have to land to be interesting; we have not put one there.

\textbf{Single-seed language-model results.} Three episodes, one seed, one run
each, on an engine version since corrected for a defect those episodes exposed.
They calibrate cost and show the interface is usable without explanation. They
do not rank models and we make no such ranking.

\textbf{The causal-confusion ablation is not run.} We designed against spurious
terrain shortcuts \citep{tien2023causal} and we show that a pure lookup table
loses money. We have \emph{not} shown that our reasoning agents are free of some
other surface shortcut. The honest test would hold physics fixed while permuting
the cosmetic identity of terrain types; an agent relying on a genuine causal
chain should be unaffected. That experiment is designed and unrun, and it is the
first thing we would do next.

\textbf{Cost accounting favours the scripted agents.} We price their CPU at a
generous \$0.10/hour and count language-model tokens at published rates. The
ratio is large enough that the accounting convention does not change the
conclusion, but it is a convention.

\section{Reproducing this}

Deterministic: a seed, a configuration and an ordered action log reproduce a
bit-identical state hash. Without that, a score difference between agents might
be luck.

\begin{lstlisting}[language=bash]
pip install -e .
wfbench bench --agents donothing,lookup,random,heuristic --seeds 1-10 --years 10
wfbench generalise --agent heuristic --train 1-8 --held-out 9-10
wfbench serve            # line-delimited JSON, for language-model agents
wfbench api              # HTTP, for web front-ends and remote agents
\end{lstlisting}

Full interface documentation, including the action schema and the three ways to
drive the engine, is in \texttt{docs/API.md}. Forty-one executable invariants
cover per-tick water mass balance and energy conservation, determinism under
action-log replay, wind-shadow directionality, turbine cut-in/cut-out, hydro
floors, and the anti-memorisation probe.

\bibliographystyle{plainnat}
\bibliography{references}

\end{document}
